\documentclass[11pt]{article}

\usepackage[margin=1in]{geometry}
\usepackage{amsmath}
\usepackage{amssymb}
\usepackage{amsthm}
\usepackage{booktabs}
\usepackage{array}
\usepackage[hidelinks]{hyperref}

\setlength{\emergencystretch}{3em}

\theoremstyle{plain}
\newtheorem{theorem}{Theorem}[section]
\newtheorem{lemma}[theorem]{Lemma}
\newtheorem{proposition}[theorem]{Proposition}
\newtheorem{corollary}[theorem]{Corollary}
\newtheorem{conjecture}{Conjecture}
\renewcommand{\theconjecture}{\arabic{conjecture}}

\theoremstyle{definition}
\newtheorem{definition}[theorem]{Definition}

\theoremstyle{remark}
\newtheorem{remark}[theorem]{Remark}

\newcommand{\R}{\mathbf{R}}
\newcommand{\Cx}{\mathbf{C}}
\newcommand{\N}{\mathbf{N}}
\newcommand{\Q}{\mathcal{Q}}
\newcommand{\re}{\operatorname{Re}}
\newcommand{\im}{\operatorname{Im}}
\newcommand{\XiH}{\Xi_N}
\newcommand{\tphi}{\widetilde{\varphi}}

\title{A counterexample to Haglund's monotonic-zeros conjecture for the Riemann
$\Xi$ approximants}
\author{Kunal Tyagi\thanks{%
\texttt{hello.jay.tyagi@gmail.com}. \emph{Use of AI.} The mathematics in this
paper --- the derivations, the computations, the verification suite, and the
text --- was produced by Claude (Anthropic) working under the author's
direction inside a structured research program; the author set the program's
objectives, adjudicated its decisions, and is responsible for the content.
Claude is a tool and is not an author. The complete working record is public:
\url{https://github.com/x67ai/riemann-rh-program}.}}
\date{October 1, 2026}

\begin{document}

\maketitle

\begin{abstract}
Haglund introduced the partial sums $\Xi_N(z) = \sum_{n \le N} \Phi_n(z)$ of
Riemann's expansion of the $\Xi$-function in incomplete gamma functions, and
conjectured that for every $N$ the zeros of $\Xi_N$ in the closed first quadrant,
listed by increasing real part, have nondecreasing imaginary parts. The
conjecture implies the Riemann hypothesis, and Haglund reported computations
indicating that a weaker form of it holds for $N \le 10$. We show that the
conjecture is false for $N = 27$: $\Xi_{27}$ has a real zero in
$(3144.8946, 3144.8947)$ and a non-real zero
$z_0 = 3143.2206824215\ldots + 0.3152587993\ldots\, i$ in the first quadrant,
whose real part is smaller. The weak form fails as well. Both facts are
certified twice, independently, in rigorous interval arithmetic --- once with Arb
ball arithmetic, once with outward-rounded interval arithmetic and proved
truncation bounds --- and $z_0$ is enclosed in a square of half-width
$10^{-48}$. The same violation occurs at $N = 24$ for the companion
approximants $\xi_N = \Xi_N + c_N$. The failure has a mechanism: on the critical
line $\Xi_N < \Xi < \xi_N$, so the real zeros of $\Xi_N$ lie in those positive
lobes of $\Xi$ that rise above a positive level, and a shallow lobe followed by
a deeper one near the height $4(N+1)^2$, where the approximant departs from
$\Xi$, leaves a non-real zero below real zeros. Two statements in Haglund's paper
are corrected: the number of real zeros of $\Xi_N$ is odd, so his table entry
$32$ at $N = 4$ is inconsistent (a numerical census finds $31$), and the
$1/x^2$ coefficient of $\Xi_N$ on the real axis is negative for every $N$. The
result refutes a sufficient condition for the Riemann hypothesis and says
nothing about the Riemann hypothesis itself: the non-real zero belongs to the
approximant, at a height where the approximant no longer tracks $\Xi$, and the
Riemann hypothesis has been verified far beyond that height.
\end{abstract}

%% SECTION 1
\section{Introduction}
\label{sec:intro}

\subsection{Haglund's approximants and his conjecture}

Write $s = \tfrac12 + iz$ and let
\begin{equation}\label{eq:Xi}
\Xi(z) = \tfrac12 s(s-1)\pi^{-s/2}\Gamma(s/2)\zeta(s) = \xi(s),
\end{equation}
Riemann's $\Xi$-function, which is even, entire, and real on the real axis; the
Riemann hypothesis is the statement that all its zeros are real. Let
$\Q = \{z : \re z \ge 0,\ \im z \ge 0\}$ be the closed first quadrant. Haglund
\cite{Hag} starts from Riemann's representation
$\Xi(z) = \int_0^\infty \cos(zt)\varphi(t)\,dt$ with
$\varphi = \sum_{n \ge 1}\varphi_n$,
$\varphi_n(t) = \exp(-n^2\pi e^{2t})(8\pi^2 n^4 e^{4.5t} - 12\pi n^2 e^{2.5t})$,
and asks what happens when $\varphi$ is replaced by the finite sum
$\sum_{n \le N}\varphi_n$. With the ``hyperbolic gamma function''
\begin{equation}\label{eq:G}
G(z;a,b) = 4\int_0^\infty \cos(2zu)\,e^{2bu - a e^{2u}}\,du
= \frac{\Gamma(b+iz,a)}{a^{b+iz}} + \frac{\Gamma(b-iz,a)}{a^{b-iz}}
\qquad (a > 0,\ b \in \R),
\end{equation}
where $\Gamma(w,a) = \int_a^\infty e^{-t}t^{w-1}\,dt$ is the upper incomplete gamma
function, this gives Riemann's series
$\Xi(z) = \sum_{n \ge 1}\Phi_n(z)$ and Haglund's \emph{$\Xi$-approximates}
\begin{equation}\label{eq:XiN}
\Xi_N(z) := \sum_{n=1}^{N}\Phi_n(z), \qquad
\Phi_n(z) := 2\pi^2 n^4\, G\bigl(\tfrac{z}{2}; n^2\pi, \tfrac94\bigr)
- 3\pi n^2\, G\bigl(\tfrac{z}{2}; n^2\pi, \tfrac54\bigr)
\end{equation}
(\cite[(1)--(3), (6), (10)--(14), pp.~1--3]{Hag}; page and equation numbers refer
to the arXiv version). Each $\Xi_N$ is even, entire, and real on the real axis.
Haglund defines: ``We say that a given function $F(z)$ has monotonic zeros in a
region $D$ of the complex plane if, when we list the zeros of $F$ in $D$ by
increasing real part, the imaginary parts of the zeros are monotone
nondecreasing'', with the side assumption ``(We assume $F$ has at most one zero
on the intersection of any vertical line with $D$.)'' \cite[p.~3]{Hag}. He then
states:

\begin{conjecture}[Haglund {\cite[p.~3]{Hag}}]\label{conj:H}
``For $N \in \N$, $\Xi_N(z)$ has monotonic zeros in $\Q$.''
\end{conjecture}

His Proposition~1 reads ``Conjecture 1 implies the Riemann Hypothesis''
\cite[pp.~3--4]{Hag}; the proof is a Hurwitz-type argument, using that
$\Xi_N \to \Xi$ uniformly on compacta and that $\Xi$ has infinitely many real
zeros. His Remark~1 records a weak form: ``A weaker form of Conjecture 1, which
still implies RH, is that there are no non-real zeros of $\Xi_N(z)$ in $\Q$ whose
real part is less than the largest real zero of $\Xi_N(z)$'', and reports that
``Computations along these lines indicate this weaker form of Conjecture 1 is
true at least for $N \le 10$'' \cite[p.~4]{Hag}, the computations being
floating-point ones for which ``no attempt has been made to establish rigorous
error bounds'' \cite[p.~3]{Hag}. The paper was published as \cite{HagJ}.

\subsection{The result}

\begin{theorem}\label{thm:main}
\begin{enumerate}
\item[(a)] $\Xi_{27}$ has a real zero $x_1 \in (3144.8946,\ 3144.8947)$, and a
second real zero $x_2 \in (3145.5998,\ 3145.5999)$.
\item[(b)] $\Xi_{27}$ has exactly one zero $z_0$, which is simple, in the open
square of half-width $3.7\cdot 10^{-11}$ about
$c = 3143.2206824215 + 0.3152587994\,i$. More precisely,
$z_0 \in z^* + [-10^{-48}, 10^{-48}]^2$ with
\[
z^* = 3143.220682421536585287281295989417800119\ldots
+ 0.315258799378214845382286373009271765\ldots\, i,
\]
where the full explicit center $z^*$ is given in Section~\ref{sec:cert}.
\item[(c)] Consequently $z_0 \in \Q$ is non-real, with $\im z_0 > 0.3152$ and
$\re z_0 < 3143.2207 < x_1$. Listed by increasing real part, the zeros of
$\Xi_{27}$ in $\Q$ do not have nondecreasing imaginary parts: Conjecture~\ref{conj:H}
is false for $N = 27$. Since $z_0$ is a non-real zero in $\Q$ whose real part is
less than the real zero $x_1$, and hence less than the largest real zero, the
weak form of Remark~1 is false for $N = 27$ as well.
\end{enumerate}
\end{theorem}

The pair $(z_0, x_1)$ violates the conjecture under every enumeration of the
zeros by increasing real part, so Haglund's side assumption (at most one zero on
each vertical line) plays no role. Statements (a) and (b) are established by two
independent computer-assisted proofs in rigorous interval arithmetic, described
in Section~\ref{sec:cert}: one in Arb ball arithmetic \cite{Arb, FLINT}, one in
outward-rounded interval arithmetic \cite{mpmath} with every truncation bounded
by a proved remainder estimate. Each was first run on a ladder of known cases
(two zeros from Haglund's table and a non-real zero of $\Xi_1$ from his
appendix).

\subsection{Scope}

Theorem~\ref{thm:main} refutes a sufficient condition for the Riemann hypothesis.
It says nothing about the Riemann hypothesis, and nothing about the zeros of
$\zeta$. The non-real zero $z_0$ is a zero of the approximant $\Xi_{27}$, not of
$\Xi$, and it lies at a height where $\Xi_{27}$ no longer tracks $\Xi$: there the
difference $\Xi - \Xi_{27}$ is of the same size as $\Xi$ itself
(Section~\ref{sec:departure}). The Riemann hypothesis has been verified
rigorously far above this height: Platt and Trudgian prove that ``all zeroes
$\beta + i\gamma$ of the Riemann zeta-function with $0 < \gamma \le 3\cdot 10^{12}$
have $\beta = 1/2$'' \cite{PT}. Nor does the paper decide Conjecture~\ref{conj:H}
for $N \le 26$: the scan that located the violation at $N = 27$
(Section~\ref{sec:departure}) is numerical, covers a window near one height for
each $N$, and is not a proof that smaller $N$ satisfy the conjecture.

\subsection{The mechanism, and two corrections to Haglund's paper}

The failure is not an accident of $N = 27$. On the critical line the approximants
sit below $\Xi$ by a positive amount: for every $N \ge 1$ and real $t$,
\[
\Xi_N(t) = \Xi(t) - Q_N(t), \qquad Q_N(t) > 0
\]
(Theorem~\ref{thm:sandwich}). Hence $\Xi_N < 0$ wherever $\Xi \le 0$, and the real
zeros of $\Xi_N$ lie in the \emph{positive lobes} of $\Xi$ (the intervals between
consecutive zeros on which $\Xi > 0$): an even number in each positive lobe
$(\gamma_j, \gamma_{j+1})$ and an odd number in the central lobe $(0, \gamma_1)$,
where $0 < \gamma_1 < \gamma_2 < \cdots$ are the positive zeros of $\Xi$. A positive lobe carries real zeros of $\Xi_N$
only if $\Xi$ rises above the level $Q_N$ there. Near the height $4(N+1)^2$ the
level $Q_N$ becomes comparable to the size of $\Xi$, and the heights of
consecutive lobes fluctuate; a lobe that stays below the level, followed by a
lobe that rises above it, leaves a pair of non-real zeros below real zeros. At
$N = 27$ the positive lobe $(3142.97945673,\ 3143.56749918)$ of $\Xi$ reaches
only $0.411$ of the level $Q_{27}$, and the next positive lobe
$(3144.49048815,\ 3146.28352838)$ reaches $1.325$ of it. The first lobe carries
no real zero of $\Xi_{27}$, and the pair $z_0, \bar z_0$ sits over it; the second
lobe carries the real zeros $x_1, x_2$ of Theorem~\ref{thm:main}. Section~\ref{sec:departure}
gives the lobe figures, which are floating-point computations, not part of the
proof.

The same sandwich corrects two statements in \cite{Hag}.
\begin{enumerate}
\item[(i)] The number of real zeros of $\Xi_N$ on $(0,\infty)$, counted with
multiplicity, is odd for every $N$ (Corollary~\ref{cor:odd}). Haglund's table
\cite[p.~4]{Hag} lists $1, 7, 15, 32, 53, 79, 113, 155, 207, 263$ real zeros for
$N = 1, \ldots, 10$; the entry $32$ at $N = 4$ is inconsistent with this, and a
numerical census of $\Xi_4$ finds $31$, the largest being
$103.3679880094135$, which is Haglund's printed value.
\item[(ii)] On the real axis $x^2\Xi_N(x)$ tends to a negative limit for every
$N \ge 1$ (Proposition~\ref{prop:tail}), whereas \cite[p.~10]{Hag} states
``Thus the coefficient of $1/x^2$ in $\Xi_N(x)$ is positive for $k \ge 3$ and
negative for $1 \le k < 3$'' (his $k$ is $N$).
\end{enumerate}

\subsection{The companion approximants, and why an ordering property}

Haglund's $\Xi_N$ differs by a constant from the more obvious truncation of
Riemann's formula,
\[
\xi_N(s) := \tfrac12 + \tfrac12 s(s-1)\sum_{n \le N} g_n(s), \qquad
\Xi_N(z) = \xi_N(\tfrac12 + iz) - c_N, \qquad
c_N = \sum_{n > N}(4\pi n^2 - 1)e^{-\pi n^2} > 0
\]
(Section~\ref{sec:riemann}). On the critical line $\xi_N$ lies \emph{above} $\Xi$
(Theorem~\ref{thm:polya}), so its real zeros lie in the negative lobes of $\Xi$,
and the same ordering property fails for it at $N = 24$
(Theorem~\ref{thm:xi24}). Neither family can satisfy the Riemann hypothesis
literally: no member of the family
$\xi_w = \tfrac12 + \tfrac12 s(s-1)\sum_n w_n g_n$ with $0 \le w_n \le 1$ and
$w \not\equiv 1$, or with $w$ finitely supported, has all its zeros on the
critical line; it has finitely many zeros there and infinitely many off it
(Theorem~\ref{thm:A}). With polynomially bounded weights, having all zeros on
the line forces $w \equiv 1$ (Theorem~\ref{thm:B}). This is why a conjecture
about approximants of this kind has to be an ordering statement such as
Conjecture~\ref{conj:H}, and not real-rootedness; and any approach through
real-rootedness in a window that grows with $N$ is a height filtration of the
Riemann hypothesis, with nothing to prove beyond the Riemann hypothesis on the
new window (Proposition~\ref{prop:F}).

\subsection{Prior work}

Haglund's conjecture was stated in 2009 \cite{Hag, HagJ} and checked by him, in
its weak form, for $N \le 10$. Ahn's thesis \cite{Ahn}, written under Haglund's
supervision, computes zeros of the component functions of the approximants
(incomplete gamma functions, hyperbolic gamma functions, and small linear
combinations of the $\Phi_n$) and restates the check: ``Computations indicate
this weaker form of Conjecture 1.6 is true at least for $N \le 10$''
\cite[p.~10]{Ahn}; it tests no $\Xi_N$ beyond $N = 10$. Haglund's Conjecture~4
concerns the zeros of the pencil $\Xi_k + t\Phi_{k+1}$, $0 \le t \le 1$
\cite[p.~11]{Hag}; its first case $k = 1$ is the subject of a 2026 preprint of
Baccaro \cite{Bac}, whose record states that ``The result covers only the first
interpolation.'' We found no test of Conjecture~\ref{conj:H} beyond $N = 10$, and
no refutation of it, in the literature we could reach; a null search is not a
proof of absence. Lagarias and Montague study a related Fourier-integral family
whose members have only finitely many zeros on the critical line; they ask
whether monotone increase of the imaginary parts of the zeros in the first
quadrant would be a sufficient condition for all zeros to be on the line, and
point to Haglund's paper as a related situation ``involving the $\xi$-function,
where such a monotonicity implies the RH'' \cite[p.~24]{LM}. Truncations that keep
all but finitely many zeros \emph{on} the line are known in another setting: Ki
proved that for $y \ge 1$ all but finitely many complex zeros of any partial sum
of the Chowla--Selberg formula are simple and on the line, and observed that the
Riemann hypothesis would follow if the words ``but finitely many'' could be
removed for infinitely many $N$ in the case $z = i$ \cite[p.~198]{Ki}. The
approximants here are the opposite case (Theorem~\ref{thm:A}). Finally,
Liu, Matiyasevich, Oesterl\'e and Zaharescu prove, for a different family of
approximants to $\xi$ built from truncated Euler products, that ``The Bounded
Riemann Hypothesis is equivalent to the Riemann Hypothesis plus the Simplicity
Conjecture'' \cite[Theorem~1.7]{LMOZ}; Proposition~\ref{prop:F} is the analogue
for the approximants of this paper.

\subsection{Organization}

Section~\ref{sec:riemann} recalls Riemann's identity and relates the two
families. Section~\ref{sec:positivity} proves the positivity and sandwich
theorems and the two corrections. Section~\ref{sec:structure} gives the
structural theorems. Section~\ref{sec:departure} describes the departure height
and the lobe scan. Section~\ref{sec:cert} is the certificate for
Theorem~\ref{thm:main}, and Section~\ref{sec:xi24} the companion violation at
$N = 24$.

\section{Riemann's identity and the two families}
\label{sec:riemann}

Let $\psi(x) = \sum_{n \ge 1} e^{-\pi n^2 x}$. Splitting
$\int_0^\infty \psi(x)x^{s/2-1}\,dx$ at $x = 1$ and using
$\psi(1/x) = \sqrt{x}\,\psi(x) + (\sqrt{x}-1)/2$ gives Riemann's identity
\cite{Rie, Edw}, valid for all $s \in \Cx$:
\begin{equation}\label{eq:riemann}
\Lambda(s) := \pi^{-s/2}\Gamma(s/2)\zeta(s) = -\frac1s - \frac{1}{1-s}
+ \sum_{n \ge 1} g_n(s), \qquad
g_n(s) = \frac{\Gamma(s/2, X)}{X^{s/2}} + \frac{\Gamma((1-s)/2, X)}{X^{(1-s)/2}},
\quad X = \pi n^2 .
\end{equation}
Hence $\xi(s) = \tfrac12 s(s-1)\Lambda(s) = \tfrac12 + \tfrac12 s(s-1)\sum_{n\ge1}
g_n(s)$. For weights $w\colon \N \to \R$ put
\begin{equation}\label{eq:xiw}
\xi_w(s) := \tfrac12 + \tfrac12 s(s-1)\sum_{n \ge 1} w_n\, g_n(s),
\end{equation}
and write $\xi_N$ for $w = \mathbf{1}_{[1,N]}$. Every $\xi_w$ satisfies
$\xi_w(1-s) = \xi_w(s)$ and $\xi_w(\bar s) = \overline{\xi_w(s)}$, so
$\xi_w(\tfrac12 + it)$ is real for real $t$; and
$|g_n(s)| \le g_n(\re s)$, since $|\Gamma(a, X)| \le \Gamma(\re a, X)$.
Throughout, $s = \tfrac12 + iz$; on the critical line $z = t$ is real, and we
write $\Xi(t) = \xi(\tfrac12 + it)$ and
$P_w(t) := \xi_w(\tfrac12 + it) - \Xi(t)$.

\begin{lemma}\label{lem:relation}
For every $n \ge 1$, $\Phi_n(z) = \tfrac12 s(s-1)g_n(s) + (4\pi n^2 - 1)e^{-\pi n^2}$.
Consequently
\[
\Xi_N(z) = \xi_N(s) - c_N, \qquad
c_N := \sum_{n > N}(4\pi n^2 - 1)e^{-\pi n^2} > 0 .
\]
\end{lemma}

\begin{proof}
Put $h(w) = X^{-w}\Gamma(w, X)$ with $X = \pi n^2$. Since $iz = s - \tfrac12$,
one has $\tfrac94 + \tfrac{iz}{2} = 2 + \tfrac s2$ and
$\tfrac94 - \tfrac{iz}{2} = 2 + \tfrac{1-s}{2}$, and the same with $1$ in place
of $2$ for $b = \tfrac54$; so definition \eqref{eq:XiN} reads
$\Phi_n = 2X^2[h(\tfrac s2 + 2) + h(\tfrac{1-s}2 + 2)]
- 3X[h(\tfrac s2 + 1) + h(\tfrac{1-s}2 + 1)]$.
The recurrence $\Gamma(w+1, X) = w\Gamma(w, X) + X^w e^{-X}$ gives
$h(w+1) = (w\,h(w) + e^{-X})/X$, and applying it twice,
$2X^2 h(a+2) - 3X h(a+1) = (2a^2 - a)h(a) + (2a - 1 + 2X)e^{-X}$. For
$a = s/2$ and for $a = (1-s)/2$ one has $2a^2 - a = \tfrac12 s(s-1)$, while
$2a - 1$ equals $s - 1$ and $-s$ respectively; adding the two gives the first
claim. Summing over $n \le N$ and using
$\sum_{n \ge 1}(4\pi n^2 - 1)e^{-\pi n^2} = \tfrac12$ (differentiate
$\theta(1/x) = \sqrt{x}\,\theta(x)$, $\theta = 1 + 2\psi$, at $x = 1$) gives the
second.
\end{proof}

Numerically $c_1 = 1.7180566258547875025\cdot 10^{-4}$,
$c_2 = 5.891258854396012012\cdot 10^{-11}$,
$c_3 = 2.9589851955455663928\cdot 10^{-20}$,
$c_4 = 2.4342009049056017\cdot 10^{-32}$ and
$c_5 = 3.443519624082649\cdot 10^{-47}$. So Haglund's zero sets are the
level sets $\{\xi_N = c_N\}$ of the truncations $\xi_N$, not their zero sets; the
two families have different zeros, and both are studied below.

%% SECTION 3
